% paper/sections/01_intro.tex
% Demonstration section showing \result, \resdelta, and \resultCI usage.
% All numeric literals in prose must resolve through citation macros.

Hallucination remains a central challenge for deployed language models.
A model that confidently asserts false information undermines the trust
required for real-world use.

We propose \textbf{HalluLens}, a probe-based detector that operates directly
on intermediate layer activations, bypassing the need for external knowledge
sources or multiple sampling passes.

\paragraph{Main result.}
Our ICR-probe achieves an AUROC of
\result{baseline_comparison}{hotpotqa:icr_probe:auroc:mean}[3]
on HotpotQA, a gain of
\resdelta{baseline_comparison}{hotpotqa:icr_probe:auroc:mean}{hotpotqa:selfcheck:auroc:mean}[3]
percentage points over SelfCheck (95\% CI
\resultCI{baseline_comparison}{hotpotqa:icr_probe:auroc:mean,hotpotqa:icr_probe:auroc_lo:mean,hotpotqa:icr_probe:auroc_hi:mean}[3]).

\paragraph{Generalization.}
The probe was trained on HotpotQA and evaluated on five additional datasets
without re-training, achieving a mean AUROC of
\result{baseline_comparison}{mmlu:icr_probe:auroc:mean}[3]
on MMLU and
\result{baseline_comparison}{sciq:icr_probe:auroc:mean}[3]
on SciQ.
The performance advantage over the SelfCheck baseline
holds across all benchmarks.

\paragraph{Contributions.}
We make the following contributions:
(1) A lightweight hallucination probe that runs at inference time on a
    Llama-3.1 or Qwen3-8B backbone with no additional sampling overhead.
(2) A reproducible benchmark suite covering six LLMsKnow tasks.
(3) Open-source code and a paper build system ensuring all reported numbers
    trace to a single CSV cell.
